\section{Pregunta N$^{\circ}$14\qquad Aldo Luna Bueno}

\begin{frame}
	\begin{enumerate}\setcounter{enumi}{13}
		\item

		      Sea $A\in\operatorname{GL}\left(\mathbb{R}^{n}\right)$ y
		      $A^{\left(k\right)}$ la matriz obtenida en la $k$-ésimo
		      paso de la eliminación gaussiana para $A$ con
		      $A^{\left(0\right)}=A$.

		      Sea
		      \begin{math}
			      A^{\left(k\right)}=
			      \left(
			      a_{r,s}^{\left(k\right)}
			      \right)
		      \end{math}
		      y
		      \begin{math}
			      a_{k}=
			      \max\limits_{r,s}
			      \left|
			      a_{r,s}^{k}
			      \right|
		      \end{math}.
		      Suponga que el pivoteo parcial es usado en la eliminación.
		      Muestre que:

		      \begin{enumerate}[a)]
			      \item

			            \begin{math}
				            \forall k\in
				            \left\{1,\dotsc,n-1\right\}:
				            a_{k}\leq
				            2^{k}a_{0}
			            \end{math}.

			      \item

			            \begin{math}
				            \forall k\in
				            \left\{1,\dotsc,n-1\right\}:
				            a_{k}\leq 2-a_{0}
			            \end{math}
			            cuando $A$ es tridiagonal.
		      \end{enumerate}
	\end{enumerate}
	\begin{solution}

		\begin{enumerate}[a)]
			\item

			      .

			\item

			      .
		\end{enumerate}

	\end{solution}
\end{frame}
